\addcontentsline{toc}{section}{Заключение}
\section*{Заключение}

В ходе курсовой работы был сформирован калькулятор «большой» конечной арифметики
\(\langle Z^8_8;\ +,\ * \rangle\) для четырех действий
\((+,\ -,\ *,\ \div)\) на основе «малой» конечной арифметики, где задано
правило «+1» и выполняются свойства коммутативности \((+,\ *)\),
ассоциативности \((+,\ *)\), дистрибутивности \(*\) относительно \(+\),
заданы аддитивная единица «a» и мультипликативная единица «b», а также
выполняется свойство: для \((\forall x)\ x * a = a.\)

\subsection*{Реализованный функционал}

В программной части была реализована «малая» конечная арифметика с действиями над однозначными элементами через диаграмму Хассе.
На основе «малой» арифметики реализована «большая» конечная арифметика с действиями над многозначными числами до восьми разрядов.  Добавлена поддержка отрицательных чисел с корректной обработкой знаков при выполнении всех действий.

Разработан консольный интерфейс для просмотра информации о программе, диаграммы Хассе.  Реализована обработка ошибок, включая деление на ноль, переполнение разрядов и некорректный формат чисел.

\subsection*{Достоинства реализации}
% Использование предвычисленных таблиц действий делает вычисления эффективными:  все действия малой арифметики вычисляются один раз при инициализации, а затем используются для быстрого доступа к значениям. Это позволяет избежать повторных вычислений при работе с многозначными числами.

Использование предвычисленных таблиц действий делает вычисления эффективными: на вход подаётся отношение порядка в виде диаграммы Хассе, на основе которой при инициализации формируются таблицы операций малой арифметики. Все необходимые значения вычисляются один раз, после чего используются для быстрого табличного доступа. Это позволяет избежать повторных вычислений при работе с многозначными числами.

\subsection*{Недостатки реализации}

Конфигурация системы задаётся в файле \texttt{config.hpp} и требует перекомпиляции при изменении. Пользователь не может менять параметры во время работы программы.

Также программа ограничена восьмью разрядами, что не позволяет работать с очень большими числами. При превышении лимита возникает ошибка переполнения. 

\subsection*{Масштабирование}

Все действий выполняются через консольный интерфейс.  Визуализация интерфейса, таблиц и диаграммы Хассе была бы удобнее  через GUI при помощи фреймворка Qt. 

Можно добавить возможность загружать конфигурацию из файла, чтобы пользователь мог менять алфавит и правило "$+1$" без повторной компиляции программы. Также можно добавить реализацию возведения числа в степень.